%% ma: L10-B12-0013
%% lop: 10
%% chuong: 5
%% bai: 12
%% dang: 10.12.4
%% loai: TL
%% muc: TH
%% dap_an: "-"
%% nguon: "chuyen-de-cac-so-dac-trung-cua-mau-so-lieu-khong-ghep-nhom-toan-10-chuong-trinh-moi.pdf (GV. Phan Nhật Linh) (P01), Bài 01 Bài tập tự luận 13"
%% kien_thuc: "Lớp 10 Bài 12 (số gần đúng, sai số tuyệt đối, độ chính xác, sai số tương đối, quy tròn số)"
%% trang_thai: chinh-thuc
%% ghi_chu: "Tự duyệt theo yêu cầu giáo viên 22/10/2026."
\begin{ex}
Ba học sinh thay nhau đo chiều cao. Bạn thứ nhất đo được $168\text{ cm}\pm1\text{ cm}$, bạn thứ hai đo được $181\text{ cm}\pm2\text{ cm}$, bạn thứ ba đo được $148\text{ cm}\pm1\text{ cm}$. Trong ba phép đo trên, phép đo nào chính xác nhất?
\loigiai{$\delta_1\le\dfrac1{168}\approx0{,}595\%$; $\delta_2\le\dfrac2{181}\approx1{,}105\%$; $\delta_3\le\dfrac1{148}\approx0{,}676\%$.

Vậy phép đo của bạn thứ nhất chính xác nhất.}
\end{ex}
